\answer{ex:06:1}{$e=(1-e^{-T_a/\tau_e})e^{-d/\tau_e}$. (a)~$0.110$ contre $4.5\cdot10^{-5}$, soit $\approx2.5\cdot10^3$ fois plus. (b)~$\tau_e^*=T_a/\ln(1+T_a/d)=0.59\,\text{s}$.}
\answer{ex:06:2}{Vitesse $\nu_x\nu_y(A_+\tau_+-A_-\tau_-)=0.8A_+$ par seconde, $\approx2900A_+$ en une heure; $\tau_-=\tau_+/0.6=33\ms$.}
\answer{ex:06:3}{$\mathbf e_1$ si $\mu^2+s_1^2>s_2^2$; ici $1.01>0.25$, et le neurone apprend $\mathbf e_1$, bien que la dispersion selon $\mathbf e_2$ soit $25$ fois plus grande: la règle trouve le vecteur principal de la matrice de corrélation, et non de covariance.}
\answer{ex:06:4}{$V=Re^{-(t_c+L-t)/\tau_\gamma}$ entre la lampe et le jus, donc $\dot V=V/\tau_\gamma$; pic d'aire $Re^{-L/\tau_\gamma}$: $0.61R$ et $0.78R$.}
\answer{ex:06:5}{Rapport signal sur bruit $1/(N+1)$, nombre d'essais $\propto N+1$: $5\cdot10^5$ essais $\approx1$~mois et $5\cdot10^{11}$ essais $\approx8\cdot10^4$~ans.}
\answer{ex:06:6}{(a)~$k_2=1/\tau_d$, $k_1=1/\tau_d-1/\tau_\gamma$. (b)~Fausse erreur $-(V/\tau_\gamma)\,\varepsilon/(1+\varepsilon)$, $\varepsilon=\tau_d/\tau_\gamma$, d'où $\tau_d\le0.105\,\text{s}$.}
\answer{ex:06:7}{$0.46$, $0.90$, $0.99$, $0.95$, et $0.70$ pour $d=3\,\text{s}$. Les points tombent sur une même courbe en fonction de la fraction de trace $F=(1-e^{-T_{\text{cue}}/\tau_e})e^{-d/\tau_e}$: apprentissage pour $F\gtrsim0.1$, choix aléatoire pour $F<10^{-2}$.}
\answer{ex:06:8}{$\eta^*=1/\bigl(2\sigma^2(N+2)\bigr)$, en $N+2$ essais; avec $m$ neurones fluctuants sur $N$, en $N(m+2)/m\ge N+2$: la parcimonie n'aide pas.}
